\section{実験}
\label{sec:experiments}

\begin{table*}[t]
\caption{\textbf{三つの OLAT ベンチマークの全 18 シーンでの比較。}テスト画像は 5,691 枚、目標画像は共通とし、元のテスト校正を使用してテスト時の当てはめは行わない。最良値は\textbf{太字}、次点は\underline{下線}で示す。反復数は RNG 30k+70k、OLAT-GS 30k+30k、GS$^3$ 100k、SSD-GS 100k（SSS は 60k）、LiSA 38k。}
\label{tab:main}
\centering\small
\setlength{\tabcolsep}{3.5pt}
\begin{tabular}{@{}l ccc ccc ccc ccc@{}}
\toprule
 & \multicolumn{3}{c}{NRHints、実撮影 (7)} & \multicolumn{3}{c}{GS$^3$、合成 (6)} & \multicolumn{3}{c}{SSS-GS、合成 (5)} & \multicolumn{3}{c}{全シーン (18)} \\
\cmidrule(lr){2-4}\cmidrule(lr){5-7}\cmidrule(lr){8-10}\cmidrule(l){11-13}
手法 & PSNR$\uparrow$ & SSIM$\uparrow$ & LPIPS$\downarrow$ & PSNR$\uparrow$ & SSIM$\uparrow$ & LPIPS$\downarrow$ & PSNR$\uparrow$ & SSIM$\uparrow$ & LPIPS$\downarrow$ & PSNR$\uparrow$ & SSIM$\uparrow$ & LPIPS$\downarrow$ \\
\midrule
RNG~\cite{fan2025rng} & 20.94 & .828 & .203 & 22.14 & .895 & .132 & 37.10 & .983 & .028 & 25.83 & .893 & .131 \\
OLAT-GS~\cite{kuang2024olat} & 27.26 & \underline{.886} & \underline{.114} & 27.11 & .940 & .079 & \textbf{40.74} & \textbf{.990} & \textbf{.016} & 30.95 & .933 & .075 \\
GS$^3$~\cite{bi2024gs3} & \underline{27.28} & .884 & .116 & \underline{31.97} & \underline{.967} & \underline{.049} & 36.81 & .983 & .028 & 31.49 & \underline{.939} & .069 \\
SSD-GS~\cite{zheng2026ssdgs} & 27.20 & .882 & .117 & 31.53 & .965 & .050 & 37.79 & \underline{.986} & \underline{.023} & \underline{31.59} & \underline{.939} & \underline{.068} \\
\midrule
\textbf{LiSA（本手法）} & \textbf{29.28} & \textbf{.909} & \textbf{.100} & \textbf{32.40} & \textbf{.970} & \textbf{.038} & \underline{40.63} & \textbf{.990} & \textbf{.016} & \textbf{33.47} & \textbf{.952} & \textbf{.056} \\
\bottomrule
\end{tabular}
\end{table*}


\subsection{実験設定}
\label{sec:setup}
\paragraph{データ。}
三つの公開 OLAT ベンチマークのすべてのシーンを使用する。NRHints~\cite{zeng2023nrhints} の七つの実撮影シーン（各シーンに学習画像 522--1500 枚、テスト画像 66--200 枚）、\gsthree~\cite{bi2024gs3} の六つの HDR 合成シーン（2000/400）、SSS-GS~\cite{dihlmann2024sss} の五つの表面下散乱シーン（500/500）である。
各ベンチマークの規約（画像は 512\,px、SSS-GS は 256\,px。HDR は白背景で $\gamma=2.2$ により表示）を維持し、すべての公式学習画像と全 5,691 枚のテスト画像を用いる。

\paragraph{ベースライン。}
点光源下の撮影画像からシーンごとに学習し、コードが公開されている近年の四つのガウシアン手法、\gsthree~\cite{bi2024gs3}、SSD-GS~\cite{zheng2026ssdgs}、OLAT Gaussians（OLAT-GS）~\cite{kuang2024olat}、RNG~\cite{fan2025rng} と比較する。公開コードと既定のスケジュールで再学習する（60k--100k 回の反復。本手法は 38k）。
すべての手法が前景マスクを目標画像の合成に用いる。OLAT-GS は表面再構成段階にも、\ours は不透明度損失にもマスクを用いる。\gsthree、SSD-GS、OLAT-GS は実撮影シーンで学習時の姿勢を微調整し、いずれの手法もテストカメラや光源を当てはめない。
RNG と OLAT-GS は複数のシーン（例えば Drums、Lego）で 21\,dB を下回る。これは、これらのシーンが手法の前提（黒背景の LDR データ、代理表面）に適合しない可能性があるためであり、得られた結果をそのまま報告する。
\todo{[E5, E6] SSS-GS~\cite{dihlmann2024sss} の安定した再現が得られた時点で比較に追加する（本研究の実行では五つのうち三つのシーンで発散した）。また、各ベースラインについて論文掲載の設定を一つ再現する。}

\paragraph{評価指標と手順。}
画像全体の PSNR、SSIM~\cite{wang2004image}、VGG-LPIPS~\cite{zhang2018unreasonable} を、まず各シーンのテストビュー間で、次にシーン間で平均する。シーンごとの差は Wilcoxon の符号付き順位検定で検証する。
すべての予測を共通の 8 ビット目標画像群で評価する（ベースラインの平均値の変化は最大 0.005\,dB）。前景マスク内の PSNR でも各ベンチマーク内の順位は変わらない（補足資料）。主結果は各手法・各シーンで一回の実行、アブレーション実験は三つの乱数シードを用いる。
開発中には \ours の旧版もこれらのベンチマークで評価した（補足資料）。最終設定は学習集合から評価用に取り分けたビューで選択し、全シーンで共通とする。

\subsection{最先端手法との比較}
\label{sec:exp_main}
\ours は 18 シーンの平均 PSNR、SSIM、LPIPS で最良の結果を得て（\cref{tab:main}）、SSD-GS を 1.89\,dB 上回る。シーンごとの比較では、すべてのベースラインに対する優位性が有意である（$p\le0.016$）。
ベンチマークをまたいで一貫して優れるベースラインはない。OLAT-GS は SSS-GS で最良であるが、\gsthree シーンでは \ours を 5.3\,dB 下回る。\gsthree は自身のシーンでは二位であるが、SSS-GS では最下位となる。
\ours はすべてのベンチマークで上位二位以内に入る唯一の手法であり、三つすべてで最良の LPIPS を得る。
ただし、優位性は一部のシーンに集中している。三つの実撮影シーンが全体の差の 46\% を占め、六つの \gsthree シーンの PSNR では \ours と \gsthree がそれぞれ三つで最良となる（LPIPS では \ours が五つで最良）。校正が正確な 11 の合成シーンでは、\ours は SSD-GS を 1.77\,dB 上回る（$p=0.014$）が、OLAT-GS に対する差は有意ではない（$p=0.12$）。

最大の改善は、移動する投影される影や半透明物体を通過する光が支配的なシーンで得られ、最良のベースラインに対して Hotdog で 1.90\,dB、Translucent で 1.91\,dB 向上する。
\cref{fig:comparison} では、花瓶の取っ手やホットドッグの影が、床や皿を表す大きなガウシアン上に投影される。プリミティブごとの可視性は影を断片化するが、アトラスを参照する画素受光点はその輪郭を再現する。
\ours は Drums（$-1.09$\,dB）と AnisoMetal（$-0.85$\,dB）では劣る。これらの細い、あるいは異方的なハイライトは、本手法の等方的ローブでは捉えられない。

実撮影シーンでは \ours が平均 2.0\,dB 上回るが、その改善はすべて Cat、Fish、Pixiu による。他の四つでは \gsthree を 0.3\,dB 下回る。
ここでの固定校正下の評価値は、外観だけでなく位置合わせにも依存する。公式姿勢には小さな誤差があり、再構成を固定する基準を設けずに学習姿勢を微調整すると、再構成がテストカメラに対して移動しうる。本手法のゲージ制約付き微調整（\cref{sec:implementation}）は、その仕組みによりこれを防ぐ。
各テストレンダリングを最適な 2D 並進で位置合わせすると（補足資料）、その影響が分かる。ベースラインは 4.0--4.6\,px の補正を要するのに対し、\ours は 1.7\,px であり、\ours の優位性は 0.7\,dB に縮小する。Cat と Pixiu では大きな差が残るが、Fish では差が消える。
\torun{E2：姿勢微調整を除いた \ours と、全手法に対する回転のみのテスト姿勢位置合わせ。}

\begin{figure*}[t]
\centering
\setlength{\tabcolsep}{0.8pt}
\renewcommand{\arraystretch}{0.5}
\newcommand{\cw}{0.131\textwidth}
\begin{tabular}{@{}ccccccc@{}}
\small テストビュー & \small 正解画像 & \small \ours（本手法） & \small SSD-GS & \small \gsthree & \small OLAT-GS & \small RNG \\[2pt]
\includegraphics[width=\cw]{figures/comparison/translucent/full.png} & \includegraphics[width=\cw]{figures/comparison/translucent/GT.png} & \includegraphics[width=\cw]{figures/comparison/translucent/LiSA-staged.png} & \includegraphics[width=\cw]{figures/comparison/translucent/SSD-GS.png} & \includegraphics[width=\cw]{figures/comparison/translucent/GS3.png} & \includegraphics[width=\cw]{figures/comparison/translucent/OLAT-Gaussians.png} & \includegraphics[width=\cw]{figures/comparison/translucent/RNG.png} \\
\footnotesize Translucent & \footnotesize PSNR & \footnotesize \textbf{34.36} & \footnotesize 32.28 & \footnotesize 32.58 & \footnotesize 29.93 & \footnotesize 27.92 \\[3pt]
\includegraphics[width=\cw]{figures/comparison/hotdog/full.png} & \includegraphics[width=\cw]{figures/comparison/hotdog/GT.png} & \includegraphics[width=\cw]{figures/comparison/hotdog/LiSA-staged.png} & \includegraphics[width=\cw]{figures/comparison/hotdog/SSD-GS.png} & \includegraphics[width=\cw]{figures/comparison/hotdog/GS3.png} & \includegraphics[width=\cw]{figures/comparison/hotdog/OLAT-Gaussians.png} & \includegraphics[width=\cw]{figures/comparison/hotdog/RNG.png} \\
\footnotesize Hotdog & \footnotesize PSNR & \footnotesize \textbf{33.97} & \footnotesize 31.93 & \footnotesize 32.08 & \footnotesize 26.85 & \footnotesize 28.49 \\[3pt]
\includegraphics[width=\cw]{figures/comparison/pikachu/full.png} & \includegraphics[width=\cw]{figures/comparison/pikachu/GT.png} & \includegraphics[width=\cw]{figures/comparison/pikachu/LiSA-staged.png} & \includegraphics[width=\cw]{figures/comparison/pikachu/SSD-GS.png} & \includegraphics[width=\cw]{figures/comparison/pikachu/GS3.png} & \includegraphics[width=\cw]{figures/comparison/pikachu/OLAT-Gaussians.png} & \includegraphics[width=\cw]{figures/comparison/pikachu/RNG.png} \\
\footnotesize Pikachu（実撮影） & \footnotesize PSNR & \footnotesize 30.44 & \footnotesize 29.58 & \footnotesize \textbf{30.71} & \footnotesize 29.46 & \footnotesize 16.31 \\[3pt]
\includegraphics[width=\cw]{figures/comparison/fail_drums/full.png} & \includegraphics[width=\cw]{figures/comparison/fail_drums/GT.png} & \includegraphics[width=\cw]{figures/comparison/fail_drums/LiSA-staged.png} & \includegraphics[width=\cw]{figures/comparison/fail_drums/SSD-GS.png} & \includegraphics[width=\cw]{figures/comparison/fail_drums/GS3.png} & \includegraphics[width=\cw]{figures/comparison/fail_drums/OLAT-Gaussians.png} & \includegraphics[width=\cw]{figures/comparison/fail_drums/RNG.png} \\
\footnotesize Drums & \footnotesize PSNR & \footnotesize 30.90 & \footnotesize 31.32 & \footnotesize \textbf{31.93} & \footnotesize 22.23 & \footnotesize 16.23 \\[3pt]
\end{tabular}
\caption{\textbf{定性的比較。}テストビューの枠内を切り出して示す（PSNR はビュー全体）。各画像は、そのシーンで最良のベースラインに対する \ours のビューごとの優位性が中央値となるビューである。上から、投影される影が支配的な二つの \gsthree シーン、各手法の差が小さい実撮影シーン、本手法が最も劣る合成シーンを示す。SSD-GS と \gsthree のプリミティブごとの影は、花瓶の取っ手とホットドッグの移動する影を断片化するが、アトラスを参照する画素受光点はこれらを再現する。一方、\ours は Drums のシンバルの細い光の筋を捉えられない。Fish を含む他のシーンは補足資料に示す。}
\label{fig:comparison}
\end{figure*}


\subsection{アトラスの寄与}
\label{sec:exp_atlas}
\Cref{tab:ablation} では、不透明、金属、毛、半透明、表面下散乱の材質を含む五つのシーン群でアトラスをアブレーションする。本手法が特に劣るシーンのうち三つを含み、完全モデルのシード間標準偏差は 0.11\,dB である。
最も細かいアトラス階層のみを使用すると 0.44\,dB、学習した可視性を \cref{eq:moment} のモーメント判定に置き換えると 0.24\,dB 低下する。ただし、これらの単一シードの差は両者の重要性を示唆するにとどまる。また、このシーン群では同一の受光点におけるアトラス可視性と、プリミティブごとの可視性または通常のシャドウマップをまだ比較しておらず、この比較は E3 で行う。

\begin{table}[t]
\caption{\textbf{アトラスのアブレーション。}五つのシーン（Drums、AnisoMetal、Translucent、FurScene、Soap）のシーン平均を示す。三つのシードを用いた行はシード間の平均と PSNR の標準偏差を報告し、単一シードの行はシード 0 の完全モデル（31.53\,dB）と比較する。$^\dagger$容量を揃えた MLP（49.1k 対 49.6k パラメータ）は受光点のコード、位置、光源・視線方向を参照するが、アトラスは参照しない。影には引き続きアトラスを用いる。}
\label{tab:ablation}
\centering\footnotesize
\setlength{\tabcolsep}{2.5pt}
\begin{tabular}{@{}lcccc@{}}
\toprule
条件 & シード数 & PSNR$\uparrow$ & SSIM$\uparrow$ & LPIPS$\downarrow$ \\
\midrule
完全モデル & 3 & 31.51{\tiny$\pm$0.11} & .954 & .047 \\
光輸送項なし & 3 & 30.15{\tiny$\pm$0.11} & .948 & .055 \\
光輸送 $\rightarrow$ 局所残差$^\dagger$ & 3 & 31.54{\tiny$\pm$0.03} & .954 & .047 \\
単一階層（$K{=}1$） & 1 & 31.09 & .952 & .049 \\
\shortstack[l]{モーメント判定のみ\\（$g\equiv 0$）} & 1 & 31.29 & .952 & .048 \\
各ガウシアンの可視性$^\ddagger$ & 3 & \TBD & \TBD & \TBD \\
ハードシャドウマップ$^\ddagger$ & 3 & \TBD & \TBD & \TBD \\
\bottomrule
\end{tabular}
\\[2pt]{\footnotesize$^\ddagger$\torun{E3（同じ受光点を使用し、光輸送にはアトラスを維持）。単一シードの行も三つのシードで実行する。}}
\end{table}

\begin{table}[t]
\caption{\textbf{光輸送項が重要となるシーン。}完全モデルと、光輸送項を除いたモデルおよび同容量の局所残差モデルとのシーン別 PSNR 差（三つのシードの対応する差の平均$\pm$標準偏差）と、完全モデルが光輸送項へ配分する線形放射輝度の割合（8 テストビューの平均）。}
\label{tab:paired}
\centering\small
\setlength{\tabcolsep}{2.5pt}
\begin{tabular}{@{}lccc@{}}
\toprule
シーン & \shortstack{光輸送なし\\との差} & \shortstack{局所残差\\との差} & 光輸送の割合 \\
\midrule
Drums & $+0.10${\tiny$\pm$0.02} & $-0.09${\tiny$\pm$0.02} & 6\% \\
FurScene & $+0.02${\tiny$\pm$0.03} & $+0.02${\tiny$\pm$0.02} & 7\% \\
AnisoMetal & $+1.42${\tiny$\pm$0.69} & $-0.18${\tiny$\pm$0.75} & 7\% \\
Translucent & $+2.29${\tiny$\pm$0.23} & $+0.41${\tiny$\pm$0.26} & 17\% \\
Soap & $+2.99${\tiny$\pm$0.30} & $-0.29${\tiny$\pm$0.11} & 58\% \\
\bottomrule
\end{tabular}
\end{table}


光輸送項の除去は、三つのシードすべてで PSNR を 1.37\,dB 低下させ、LPIPS を 0.047 から 0.055 へ悪化させる。この低下は材質に依存する（\cref{tab:paired}）。Drums と FurScene では無視できる程度、AnisoMetal では 1.42\,dB、Translucent と Soap ではそれぞれ 2.29、2.99\,dB である。
モデルは最初の三つのシーンでは線形放射輝度の 6--7\% を光輸送項に配分するが、Translucent では 17\%、Soap では 58\% を配分する。全 18 シーンでは、不透明な合成物体で 6--9\%、表面下散乱物体で 20--58\% となる。これは学習された配分であり、物理的な分解ではない。

アトラスの効果と追加の放射輝度経路の効果を分けるため、アトラスによる影を維持しつつ、アトラスの光輸送を同容量の MLP に置き換える。この MLP は受光点のコード、位置、光源方向、視線方向を参照するが、アトラスは参照しない。
この局所残差は平均で完全モデルと同等である（31.54 対 31.51\,dB）。したがって、これらの材質の大半では、光輸送項の主な役割は影の影響を受けない非負の放射輝度経路を提供することであり、局所ネットワークでも同様に実現できる。
アトラスが上回るのは、光が可視表面に到達する前に物体を通過する様子が明確に見える Translucent のみである（$+0.41\pm0.26$\,dB。三つのシードすべてで正の差だが、有意ではない）。一方、局所残差は Soap（$-0.29\pm0.11$\,dB）と、わずかに Drums で上回る。AnisoMetal では完全モデルのシード間変動も大きい（26.4--27.9\,dB）。
したがって、本研究の全体的な優位性を光輸送項に帰することはせず、ライト空間の可視性と最適化スケジュールが主因であるという仮説を立てる。これは E3 と E4 で検証する。
光源に対する相対量を用いる設計は、光源の外挿で最も重要になると考えられるが、既存ベンチマークはそれを検証していない。テスト光源と最近傍の学習光源との角度差の中央値は $1.1^\circ$ にすぎない。
\torun{E1（\cref{tab:holdout}）：アブレーション用シーン群で三つのシードを用いる。「アトラスは局所残差より劣化が小さい」という主張を前提とせず、検証する。}

\begin{table}[t]
\caption{\textbf{学習から除外した光源方向}（\torun{E1}）。五つのアブレーション用シーンにおける、公式テストビューと、光源方向が評価用に取り分けられた 30$^\circ$ 角度区間に入る学習集合内のビューの平均 PSNR。各手法をそれらのビューを除いて再学習する。$^\dagger$三つのシードを使用。}
\label{tab:holdout}
\centering\small
\begin{tabular}{@{}lccc@{}}
\toprule
手法 & 公式テスト & 除外した光源 & 低下幅 \\
\midrule
\gsthree & 31.25 & \TBD & \TBD \\
SSD-GS & 31.01 & \TBD & \TBD \\
OLAT-GS & 28.67 & \TBD & \TBD \\
\ours、局所残差$^\dagger$ & 31.54 & \TBD & \TBD \\
\ours$^\dagger$ & 31.51 & \TBD & \TBD \\
\bottomrule
\end{tabular}
\end{table}

\begin{figure}[t]
\centering
\setlength{\tabcolsep}{0.8pt}
\renewcommand{\arraystretch}{0.5}
\newcommand{\lw}{0.322\linewidth}
\begin{tabular}{@{}ccc@{}}
\small 表示領域の損失 & \small 放射量による段階化 & \small 正解画像 \\[2pt]
\includegraphics[width=\lw]{figures/loss_domain/lego/control_crop.png} & \includegraphics[width=\lw]{figures/loss_domain/lego/ours_crop.png} & \includegraphics[width=\lw]{figures/loss_domain/lego/gt_crop.png} \\
\footnotesize 23.02\,dB & \footnotesize \textbf{30.25\,dB} & \footnotesize Lego \\[3pt]
\includegraphics[width=\lw]{figures/loss_domain/drums/control_crop.png} & \includegraphics[width=\lw]{figures/loss_domain/drums/ours_crop.png} & \includegraphics[width=\lw]{figures/loss_domain/drums/gt_crop.png} \\
\footnotesize 24.10\,dB & \footnotesize \textbf{30.10\,dB} & \footnotesize Drums \\[3pt]
\end{tabular}
\caption{\textbf{損失領域と幾何形状。}二つの実行で初期化、スケジュール、乱数状態、その他すべての設定を共有し、損失を表示用に符号化した放射輝度上で計算するか、目標値を徐々に移行するウォームアップの後に線形放射輝度上で計算するかだけを変える。16k 回の反復を行い、学習集合から評価用に取り分けた 200 ビューの PSNR を示す。両シーンですべての 200 ビューが改善する。}
\label{fig:loss_domain}
\end{figure}


\subsection{最適化の寄与}
\label{sec:exp_optimization}
\paragraph{損失領域。}
最初の実行で失敗した Lego と Drums において、16k 回の反復からなる二つの実行で初期化、ウォームアップ、スケジュール、乱数状態を共有し、損失だけを変える。対照条件は全期間にわたり表示用に符号化した放射輝度を用い、本手法は段階 I と II に従う。
学習集合から評価用に取り分けた 200 ビューで、表示領域の損失による幾何形状は崩れる（\cref{fig:loss_domain}）。一方、放射量に基づくスケジュールは PSNR を Lego で 23.02 から 30.25\,dB、Drums で 24.10 から 30.10\,dB に向上させる（LPIPS は 0.122$\to$0.057、0.060$\to$0.034）。すべてのビューで改善し、改善幅の中央値はそれぞれ 7.2、6.0\,dB である。
この実験の範囲は、二つの HDR シーン、単一シード、16k 回の反復に限られる。実撮影シーンと SSS-GS シーンでは、他の点も異なる \ours の旧来の表示領域版が 0.2--0.4\,dB 高い結果を得る（補足資料）。
\torun{E4：アブレーション用シーン群で三つのシードを用い、(a) 全期間で表示領域の損失、(b) ウォームアップなし、(c) 段階 II で画素受光点、(d) 段階 III なし、(e) $\Gamma$ 内へのオフセットを検証する。}

\paragraph{受光点と微調整。}
Lego、Drums、Soap の 30k 時点の幾何形状を固定した実験（補足資料の表 S6）では、線形放射輝度ではなく表示領域で微調整すると、両種の受光点で三指標すべてが改善する（PSNR はガウシアンで $+0.20$\,dB、画素で $+0.28$\,dB）。
画素受光点は最良の SSIM と LPIPS（0.976、0.029）を得る一方、ガウシアン受光点の PSNR は 0.19\,dB 高い。本研究では知覚品質を重視して画素を採用する。

\subsection{計算コスト}
\label{sec:exp_cost}
Drums、FurScene、Soap において他の処理がない GPU で計測すると（補足資料の表 S5）、\ours は \gsthree と SSD-GS より 1.6--5.5 倍速く学習する。Drums のガウシアン数は \ours が 228k であるのに対し、両ベースラインは約 550k である。全 18 シーンでは計測条件がそれほど統一されていないが、平均学習時間は本手法が 21 分、ベースラインが 55--107 分である。
毎フレームで新しい光源を用いても、光源パスを含むレンダリング速度は毎秒 63--104 フレームであり、対話的な操作が可能である。ただし、光源パスのため、小さいシーンでは \gsthree と SSD-GS より遅い。

\subsection{限界}
\label{sec:limitations}
等方的ローブはヘアライン加工された金属や細い光の筋を捉えられない。Drums では、正解画像で細いハイライトとして抽出した 0.27\% の画素が、\gsthree に対する本手法の超過二乗誤差の 72\% を占める。
モーメント判定は複数の表面が混在する集約範囲に対して弱い事前知識にとどまり、中心深度を用いると大きな斜めのガウシアンの深度位置がずれる。チェビシェフ判定や四モーメント判定~\cite{donnelly2006variance,peters2015moment}、交差点の深度の利用は自然な改善案である。
一枚の透視投影アトラスは光源が物体の外部にあることを仮定する。光輸送項は立体角係数と入射点での距離減衰を無視し、光源から見えない領域だけを経由して受光点に到達する光は、一回のライト空間での集約では扱えない。
